\newcommand{\SPEightWorlds}{900}
\newcommand{\SPEightRecords}{4500}
\newcommand{\SPEightExecutedRuns}{4050}
\newcommand{\SPEightCertified}{450}
\newcommand{\SPEightProposedSuccess}{0.331}
\newcommand{\SPEightDegradedSuccess}{0.290}
\newcommand{\SPEightNaiveSuccess}{0.288}
\newcommand{\SPEightPerfectSuccess}{0.456}
\newcommand{\SPEightHOneEffect}{0.054}
\newcommand{\SPEightHTwoEffect}{1.561}
\newcommand{\SPEightHThreeEffect}{261.3}
\newcommand{\SPEightMaxScale}{64}
\newcommand{\SPEightMaxRobots}{128}


\subsection{SP8: escala y coordinación bajo red imperfecta}
\label{subsec:sp8}

Cuando dos componentes de la red dejan de observar una interacción común, cada uno puede alcanzar un Nash visible y conservar un conflicto global. SP8 mide este fallo al escalar las coaliciones discretas de SP7 y variar radio, retardo y pérdida. El número de robots es $N=2A$ y la campaña no incorpora una planta física.

\subsubsection{Escala y observabilidad}

El retardo produce copias obsoletas, la pérdida deja vistas incompletas y una partición oculta por completo los conflictos entre componentes (Figura~\ref{fig:sp8-scenario}).

\begin{figure}[!htbp]
  \centering
  \begin{spdiagram}
    \spdrawarena{SP8 \textbar\ escala, vecindad y pérdida de información}
    \node[sprobot] (c1) at (1.25, 3.55) {$C_1$};
    \node[sprobotgreen] (c2) at (3.65, 3.55) {$C_2$};
    \node[sprobotpurple] (c3) at (2.45, 1.65) {$C_3$};
    \draw[spcomm] (c1)--(c2)--(c3)--(c1);
    \node[spannotation] at (2.45, 4.35) {grafo nominal};

    \node[sprobot] (d1) at (6.20, 3.55) {$C_1$};
    \node[sprobotgreen] (d2) at (8.60, 3.55) {$C_2$};
    \node[sprobotpurple] (d3) at (7.40, 1.65) {$C_3$};
    \draw[spcommdelayed] (d1)--(d2);
    \draw[draw=SPRed,densely dotted,line width=0.9pt] (d2)--(d3);
    \node[spannotation,text=SPRed] at (7.95, 2.55) {pérdida};
    \node[spannotation] at (7.40, 4.35) {retardo y paquetes};

    \node[sprobot] (p1) at (10.40, 3.55) {$C_1$};
    \node[sprobotgreen] (p2) at (12.55, 1.65) {$C_2$};
    \draw[draw=SPRed,dashed,line width=0.9pt] (p1)--(p2);
    \node[spannotation] at (11.48, 4.35) {partición};
    \node[spannotation,text=SPRed] at (11.48, 2.55) {conflicto no visible};

    \draw[spcomm] (0.40, 0.42)--(1.20, 0.42);
    \node[splegend,anchor=west] at (1.32, 0.42) {mensaje recibido};
    \draw[spcommdelayed] (4.40, 0.42)--(5.20, 0.42);
    \node[splegend,anchor=west] at (5.32, 0.42) {mensaje retardado};
    \draw[draw=SPRed,densely dotted,line width=0.9pt] (8.55, 0.42)--(9.35, 0.42);
    \node[splegend,anchor=west] at (9.47, 0.42) {información ausente};
  \end{spdiagram}
  \caption{Incremento de SP8: el juego regula conflictos visibles, mientras el retardo, la pérdida y la partición alteran la información disponible.}
  \label{fig:sp8-scenario}
  \viuownsource
\end{figure}

\subsubsection{Problema formal de referencia}

Sean $\mathcal A_8$ las coaliciones, $\mathcal R_i^8$ sus rutas y $\mathcal E_{ir}^8$ los recursos. $b_{ir}^8$ es coste normalizado y $P_8(\boldsymbol r)$ cuenta recursos compartidos. El oráculo resuelve
\begin{equation}
\label{eq:sp8-global-oracle}
 \min_{\boldsymbol r\in\prod_i\mathcal R_i^8}
 J_8(\boldsymbol r)
 =\sum_{i\in\mathcal A_8} b_{i r_i}^8+\lambda_8P_8(\boldsymbol r).
\end{equation}
Con dos rutas se enumeran $2^A$ perfiles, hasta 4096. Fuera de ese dominio no se informa brecha.

\subsubsection{Métodos de comparación}

La retransmisión periódica intenta cerrar vistas mediante anuncios repetidos; la ablación transmite únicamente al cambiar de ruta. La vista perfecta aísla el efecto de la red y~\eqref{eq:sp8-global-oracle} certifica el óptimo en los tamaños enumerables. La Tabla~\ref{tab:sp8-methods} sitúa estos tratamientos frente a los antecedentes de optimización asíncrona y consenso.

\begin{table}[!htbp]
  \centering
  \caption{Métodos y literatura comparada para SP8.}
  \label{tab:sp8-methods}
  \viutablefont
  \setlength{\tabcolsep}{1.3pt}
  \begin{tabularx}{\textwidth}{p{2.0cm} p{1.45cm} p{2.25cm} p{2.0cm} p{1.75cm} Y}
    \toprule
    Método/fuente & Clase & Coste & Arquitectura & Parámetros & Garantía/límite \\
    \midrule
    Gradiente asíncrono \parencite{tsitsiklis1986asynchronousOptimization} & Convexo, bajo los supuestos citados & $O(T(c_g+d_i d))$ por agente & Distribuido; valores obsoletos & Paso, retardo, tolerancia & Convergencia convexa. No cubre el juego discreto. \\
    Consenso con conmutación/retardo \parencite{olfatiSaber2004switchingConsensus} & Sistema lineal & $O(T|E_C|d)$ global & Distribuido; grafos delimitados & Pesos, $\Delta t$, $\tau_d$ & Consenso bajo conectividad. No decide rutas. \\
    ACBBA con pérdidas \parencite{rantanen2018lossyAcbba} & MRTA; heurística & $O(T(KB_i+d_iK))$ y mensajes & Asíncrona entre pares & Paquete, marcas de tiempo, pérdida & Resultados para tareas exclusivas; modelo distinto. \\
    Oráculo exhaustivo & Catálogo finito & $O(2^A A^2)$; memoria $O(A)$ & Central; información global & Límite de 4096 perfiles & Óptimo certificado solo dentro del límite. \\
    Respuesta local con versiones & Juego finito & $O(A d_i)$ por barrido; $32$ bytes/mensaje & Vecinal; rutas recibidas & Radio, retardo, pérdida, frecuencia & Nash visible; persisten conflictos no observados. \\
    \bottomrule
  \end{tabularx}
  \viusource{elaboración propia a partir de las referencias citadas}
\end{table}

\subsubsection{Juego sobre información visible}

Sea $G_C^8=(\mathcal A_8,\mathcal E_C^8)$ un grafo no dirigido y estático durante una ejecución. El número de conflictos visibles es
\[
 P_8^G(\boldsymbol r)=
 \sum_{\{i,j\}\in\mathcal E_C^8}
 |\mathcal E_{i r_i}^8\cap\mathcal E_{j r_j}^8|.
\]
Cada coalición usa la última ruta y versión recibidas y evalúa
\begin{equation}
 U_i^8(r_i,\boldsymbol r_{-i})=-b_{i r_i}^8
 -\lambda_8\sum_{j\in\mathcal N_i^8}
 |\mathcal E_{i r_i}^8\cap\mathcal E_{j r_j}^8|.
 \label{eq:sp8-visible-utility}
\end{equation}

\begin{contribucion}{potencial de red y frontera de observabilidad}
La función
\[
 \Phi_8^G(\boldsymbol r)=-\sum_i b_{i r_i}^8-\lambda_8P_8^G(\boldsymbol r)
\]
es un potencial exacto de~\eqref{eq:sp8-visible-utility}. Toda secuencia asíncrona de mejoras estrictas termina en un Nash visible tras, como máximo, $\prod_i|\mathcal R_i^8|-1$ cambios aceptados.

La identidad y la terminación finita exigen coherencia: en cada actualización aceptada, las rutas vecinas usadas por el agente deben coincidir con el perfil visible vigente. El retardo o la pérdida pueden romper esa condición, de modo que la trayectoria observada se evalúa empíricamente en lugar de tratarse como una secuencia de mejoras de $\Phi_8^G$.

\begin{proposicion}[Límite informativo bajo partición permanente]
\label{prop:sp8-partition-limit}
Considérese una familia de juegos de rutas cuyo coste global contiene interacciones entre coaliciones situadas en componentes distintas de $G_C^8$ y no incluidas en sus observaciones locales. Ningún algoritmo determinista que utilice exclusivamente historiales intracomponente puede garantizar para toda la familia una solución globalmente óptima ni detectar todas las interacciones remotas no observables.
\end{proposicion}
Una partición no siempre colisiona: prioridades o mapa previo resuelven familias particulares. La imposibilidad exige acoplamiento remoto no observable (Anexo~\ref{app:sp8-proofs}).
\end{contribucion}

Con $s$ copias y pérdidas independientes de probabilidad $p_{\mathrm{loss}}$, la probabilidad de perderlas todas es $p_{\mathrm{loss}}^s$. Las pérdidas correlacionadas y las particiones requieren modelos distintos.

\subsubsection{Protocolo de comunicación}

Un cambio de ruta genera anuncios de 32 bytes con identificador, versión y ruta. Las activaciones asíncronas aplican retardo entero y pérdida Bernoulli; el tratamiento periódico retransmite el estado, mientras la variante por evento envía únicamente los cambios.

Las rutas abstraen los cuerpos de SP3--SP4. Por ello, los conflictos y el Nash pertenecen a la capa estratégica y no describen contacto, seguimiento o barreras físicas.

\subsubsection{Diseño experimental}

Treinta semillas producen \SPEightWorlds{} mundos y \SPEightExecutedRuns{} ejecuciones. La unidad independiente es la combinación tamaño--semilla--régimen; los 720 pares se reducen a 180 instancias para la inferencia.

\begin{table}[!htbp]
  \centering
  \caption{Diseño experimental de SP8.}
  \label{tab:sp8-design}
  \viutablefont
  \begin{tabularx}{0.94\textwidth}{>{\raggedright\arraybackslash}p{2.4cm} Y}
    \toprule
    Elemento & Especificación \\
    \midrule
    Escala & $A\in\{4,8,12,16,32,64\}$ coaliciones y $N=2A$ robots. \\
    Red & Nominal, radio limitado, retardo, pérdida $0.25$ y régimen combinado con pérdida $0.35$. \\
    Contrastes & Periódico frente a evento para retransmisión; periódico frente a vista perfecta para observabilidad; oráculo en \SPEightCertified{} mundos con $A\leq12$; perfil aleatorio como referencia inferior. \\
    Métricas & Conflictos, tasa libre de conflictos, mensajes, bytes y CPU. \\
    Inferencia & Promedio por instancia, IC \mbox{bootstrap}, Wilcoxon pareado y corrección de Holm. Los 450 registros no enumerados del oráculo permanecen marcados como no ejecutados, para \SPEightRecords{} filas totales. \\
    \bottomrule
  \end{tabularx}
  \viusource{elaboración propia a partir del protocolo experimental de SP8}
\end{table}

\subsubsection{Resultados}

La retransmisión periódica alcanzó una tasa libre de conflictos de \SPEightProposedSuccess{}, frente a \SPEightNaiveSuccess{} para la variante por evento y \SPEightPerfectSuccess{} con información perfecta. En red nominal, su tasa descendió de $0.933$ con cuatro coaliciones a $0.667$ con doce, $0.433$ con dieciséis y cero con sesenta y cuatro (Tabla~\ref{tab:sp8-results}).
Para cada $A$, el generador incluye al menos dos perfiles sin conflicto: todas las coaliciones por la ruta 0 o todas por la ruta 1. El valor cero en $A=64$ corresponde a un fallo del método sobre un catálogo factible. La brecha frente al oráculo permanece sin certificar cuando $A>12$.

\begin{table}[!htbp]
  \centering
  \caption{Calidad y coste medios de los métodos ejecutados en SP8.}
  \label{tab:sp8-results}
  \resizebox{0.94\textwidth}{!}{\begin{tabular}{lrrrr}
\toprule
Método & Libre de conflicto & Pares & Mens./agente & CPU [ms] \\
\midrule
Local periódico versionado & 0.331 & 2.37 & 316.4 & 40.75 \\
Local solo por evento & 0.288 & 2.63 & 22.7 & 11.25 \\
Mejor respuesta con información perfecta & 0.456 & 1.11 & 0.0 & 11.32 \\
Perfil inicial aleatorio & 0.056 & 5.24 & 0.0 & 0.00 \\
\bottomrule
\end{tabular}
}
  \viuownsource
\end{table}


En los mundos degradados, la política periódica elevó en $0.054$ la tasa libre de conflictos frente a la política por evento (IC del 95\,\% $[0.036;0.072]$, $p_{\mathrm{Holm}}=1.45\times10^{-7}$) y añadió $261.3$ mensajes por agente (IC $[225.3;298.6]$). Frente a información perfecta, conservó $1.561$ pares de conflicto adicionales por instancia y régimen (IC $[1.336;1.813]$). Calidad y tráfico se reportan como estimandos separados.


En los 450 mundos con enumeración completa, la brecha media normalizada frente al oráculo fue $0.497$, calculada como $(J_8-J_8^\star)/\max\{|J_8^\star|,10^{-12}\}$. Los costes base hicieron $J_8^\star>0$ en todas esas instancias, por lo que la guarda numérica permaneció inactiva. La estimación cubre únicamente los tamaños enumerados.

Los regímenes de radio y combinado produjeron $2.25$ componentes en promedio y una fracción conexa de $0.389$; en los demás grafos, $\lambda_2(L)=0$.
La política periódica dejó $4.133/4.211$ conflictos remotos y cerró las vistas y el rezago.
La política por evento conservó rezagos $0.277/0.367$ y fracciones de vistas ausentes $0.157/0.207$, con un máximo de cuatro versiones. Estas magnitudes cuentan versiones simuladas y no tiempo o paquetes físicos.

La campaña alcanzó 128 robots, donde la tasa libre de conflictos nominal fue cero. El oráculo certificó hasta doce coaliciones y los tiempos de CPU describen únicamente la implementación Python evaluada.


La campaña caracteriza observabilidad, retransmisión y conflicto. Evaluar escalabilidad operacional exigiría además memoria, energía, hardware y una planta.

\subsubsection{Síntesis}

La retransmisión periódica recuperó $0.054$ de tasa libre de conflictos y añadió $261.3$ mensajes por agente. La coordinación se degradó con la escala hasta alcanzar tasa cero en $A=64$, y la brecha media fue $0.497$ en el dominio certificado por el oráculo. El potencial garantiza terminación con copias coherentes; bajo vistas obsoletas, la trayectoria tiene respaldo empírico. Estas cifras cuantifican el coste de paliar pérdidas de información en el ejecutor discreto.
